
\documentclass{article}
\usepackage{colortbl}
\usepackage{makecell}
\usepackage{multirow}
\usepackage{supertabular}

\begin{document}

\newcounter{utterance}

\centering \large Interaction Transcript for game `codenames', experiment `ambiguous', episode 2 with reflexion{-}llama{-}3.1{-}8b{-}instruct.
\vspace{24pt}

{ \footnotesize  \setcounter{utterance}{1}
\setlength{\tabcolsep}{0pt}
\begin{supertabular}{c@{$\;$}|p{.15\linewidth}@{}p{.15\linewidth}p{.15\linewidth}p{.15\linewidth}p{.15\linewidth}p{.15\linewidth}}
    \# & $\;$A & \multicolumn{4}{c}{Game Master} & $\;\:$B\\
    \hline

    \theutterance \stepcounter{utterance}  
    & & & \multicolumn{2}{p{0.3\linewidth}}{
        \cellcolor[rgb]{0.9,0.9,0.9}{
            \makecell[{{p{\linewidth}}}]{
                \texttt{\tiny{[GM$|$GM]}}
                {'hidden': {'team': ['poop', 'beetle', 'squash', 'peep', 'pink', 'meter', 'rush', 'plot', 'rocket'], 'innocent': ['corn', 'scale', 'lean', 'nag', 'stem', 'box', 'clip'], 'opponent': ['pump', 'gag', 'flush', 'bowl', 'stick', 'pawn', 'sap', 'repair'], 'assassin': []}, 'revealed': {'team': {'team': [], 'innocent': [], 'opponent': [], 'assassin': []}, 'opponent': {'team': [], 'innocent': [], 'opponent': [], 'assassin': []}}}
            }
        }
    }
    & & \\ \\

    \theutterance \stepcounter{utterance}  
    & & \multicolumn{4}{p{0.6\linewidth}}{
        \cellcolor[rgb]{0.9,0.9,0.9}{
            \makecell[{{p{\linewidth}}}]{
                \texttt{\tiny{[P1$\langle$GM]}}
                \texttt{You have attempted this same task before and did not succeed. These are the plans you wrote after those attempts.} \\
\texttt{Trial \#1: New plan: I got stuck because my clue was too narrow and only related to a few words. What is worse, I didn't consider the distractor words and the kill words, which could have given me a hint about the types of words that are not related to my clue. Next time, I will choose a clue that relates to a larger subset of words, including some of my team's words and some distractor words. I will also make sure to avoid the kill words. For example, I could choose a clue like "plant" that relates to words like "poop", "beetle", "squash", "plot", and "rocket", as well as distractor words like "corn" and "stem".} \\
\\ 
\\ 
\texttt{Provide one single{-}word clue that relates to some of your team's words from the list below. You can choose to how many and to which words your clue relates to.} \\
\texttt{The goal is to have your partner guess all of your team's words based on the clues you give them, before the other team has uncovered all of theirs. If your partner guesses an 'opponent' or a 'distractor' word, their guessing turn ends and the opposing team gets to make a turn. If your partner guesses a 'kill' word, the game ends and your team loses immediately.} \\
\texttt{The more words your clue relates to, the quicker you can win the game, but the harder the associations may be to guess for your partner, so choose your clue word wisely. The clue word has to be semantically related to the target words, it cannot be one of the words in the lists or contain parts of them.} \\
\\ 
\texttt{Always give your single{-}word clue and your comma{-}separated list of related target words in the following format and make your answers as short as possible, never include any other text than is required in this form:} \\
\\ 
\texttt{CLUE: <WORD>} \\
\texttt{TARGETS: <TARGETS>} \\
\\ 
\texttt{Your team words are: poop, beetle, squash, peep, pink, meter, rush, plot, rocket.} \\
\texttt{Your opponent's team words are: pump, gag, flush, bowl, stick, pawn, sap, repair.} \\
\texttt{Distractor words are: corn, scale, lean, nag, stem, box, clip.} \\
\texttt{The kill words are: .} \\
            }
        }
    }
    & \\ \\

    \theutterance \stepcounter{utterance}  
    & \multicolumn{4}{p{0.6\linewidth}}{
        \cellcolor[rgb]{0.9,0.9,0.9}{
            \makecell[{{p{\linewidth}}}]{
                \texttt{\tiny{[P1$\rangle$GM]}}
                \texttt{CLUE: plot} \\
\texttt{TARGETS: poop, beetle, squash, plot, rocket} \\
            }
        }
    }
    & & \\ \\

    \theutterance \stepcounter{utterance}  
    & & & \multicolumn{2}{p{0.3\linewidth}}{
        \cellcolor[rgb]{0.9,0.9,0.9}{
            \makecell[{{p{\linewidth}}}]{
                \texttt{\tiny{[GM$|$GM]}}
                {'player': 'Cluegiver', 'type': 'clue is morphologically related to word on the board', 'utterance': 'CLUE: plot\nTARGETS: poop, beetle, squash, plot, rocket', 'message': "Your clue 'plot' is morphologically similar to the word plot, please choose another clue word.", 'clue': 'plot', 'similar_board_word': 'plot'}
            }
        }
    }
    & & \\ \\

    \theutterance \stepcounter{utterance}  
    & \multicolumn{4}{p{0.6\linewidth}}{
        \cellcolor[rgb]{0.95,0.95,0.95}{
            \makecell[{{p{\linewidth}}}]{
                \texttt{\tiny{[P2$\rightarrow$P1]}}
                \texttt{Post{-}episode Reflexion (not sent to GM):} \\
\texttt{New plan: I got stuck because my clue was too narrow and only related to a few words. What is worse, I didn't consider the distractor words and the kill words, which could have given me a hint about the types of words that are not related to my clue. Next time, I will choose a clue that relates to a larger subset of words, including some of my team's words and some distractor words. I will also make sure to avoid the kill words. For example, I could choose a clue like "plant" that relates to words like "poop", "beetle", "squash", "plot", and "rocket", as well as distractor words like "corn" and "stem". However, I also need to make sure that my clue does not relate to any of the opponent's team words or the kill words. To achieve this, I will carefully examine the lists and choose a word that has a broad semantic connection to the target words, but is not too specific or too general.} \\
            }
        }
    }
    & & \\ \\

\end{supertabular}
}

\end{document}
